\cotefile{70}
\scanpage{111}
% archive sheet number: 110
Icosaèdre géométrique et représentations du groupe $\Gamma$

Soit $\Pi=(S,A,F;r)$ l'icosaèdre combinatoire-type, épinglé par le \unsure{repère|racine} $r=\deleted{\unsure{(s_0)}}(s_0,\unsure{\mathrm{ar}|\mathrm{dr}}(s_0,s_1),\unsure{\mathrm{fl}}(s_0,s_1,s_2))$, $\Gamma$ son \unsure{gpe|groupe} des autom. Une \uline{représentation}\inserted{ou \uline{réalisation}} \uline{affine} \deleted{\unsure{régulière}} de $\Pi$ dans un \deleted{\unsure{espace}} \inserted{fibré} affine $E$ (sur schéma $S$) de dim. 3 est la donnée d'applications
\[
\begin{aligned}
 S&\xrightarrow{\unsure{\rho_S|\varphi_S}}\text{ens des sections de }E\text{ sur }S\\
 A&\xrightarrow{\rho_A}\text{ens des \inserted{ss} fibrés en droites affines de }E\\
 F&\xrightarrow{\rho_F}\text{ens des sous-fibrés en plans affines}
\end{aligned}
\]
respectant les relations d'incidence, \deleted{et telles que} On dit qu'elle est \deleted{\unsure{non dégénérée}} \deleted{\unsure{quasi ``acceptable''}} ``acceptable'' si \deleted{\unsure{les}} les sommets d'une face ont des images \deleted{distinctes} \inserted{non alignées} --- ce qui implique que $\rho_A,\rho_F$ sont \unsure{connus|connues} quand on connaît $\rho_S$, \deleted{on dit que $\rho$ est fidèle si \unsure{deux} figures \deleted{\unsure{de dim.} $\unsure{t|1}$} de $\Pi$ sont incidentes \inserted{\unsure{si et seulement si}} \unsure{elles ont leurs} \unsure{images} réalisations qui le sont} \unsure{qui peut être donnée} \uline{arbitrairement}.

On dit que $\rho$ est \uline{fidèle}, \deleted{si elle est \unsure{injective}} \deleted{elle est acceptable, et} si \deleted{sur $S,A,F$} deux figures dont les réalisations \inserted{sur une fibre} sont incidentes sont incidentes. \inserted{(Il suffit de l'exiger dans le cas des sommets et arêtes, et sommets et faces)} Cela implique qu'elle est acceptable (car si \inserted{les trois} $s,s',s''$ d'une face \deleted{\unsure{sont}} ont images alignées, $\rho(s)$ serait incident à $\rho(\deleted{\unsure{s}}\{s',s''\})$ sans que $s$ le soit à $\{s',s''\}$), et $\rho_S,\rho_A,\rho_F$ sont injectives (si p.ex. $s\ne s'$ et $a$ est une arête \deleted{\unsure{par}} à laquelle $s$ et non $s'$ sont incidents, si on avait $\rho(s)=\rho(s')$ on aurait $\rho(s')$ incident à $\rho(a)$ et $s'$ \unsure{sera} inc. à $a\unsure{\ldots}$)
